\section{Conclusion}

This thesis proposes a survey experiment to examine how individuals in Germany form beliefs about the labor-market 
and social integration of Ukrainian and Syrian migrants. By independently varying migrants' country of origin, 
reason for migration, human capital, German-language proficiency, and gender, the factorial vignette design makes it 
possible to identify how these characteristics affect beliefs about employment, income, discrimination, and social 
integration. In particular, the design allows the study to examine whether differences between Ukrainian and Syrian 
migrants depend on refugee status and human capital.

The thesis considers how respondents' beliefs are related to their personal contact experiences and the information 
available to them. Measures of personal contact and the quality of contact, media exposure and tone, and memorable experiences and 
media information provide a way to examine whether differences in respondents' information environments are associated 
with their evaluations of migrant integration. The design therefore brings together two complementary perspectives: 
the characteristics of migrants that shape how they are perceived and the experiences and information through which 
respondents may form and update their beliefs. Additional measures of beliefs about migrant communities and integration 
over time provide further context for understanding these patterns. Importantly, the design distinguishes between the 
causal effects of experimentally varied migrant characteristics and the associational relationships between respondents' 
beliefs, contact experiences, and information exposure, providing a clear framework for interpreting the resulting evidence.

More broadly, understanding beliefs about migrant integration is important because public perceptions need not 
correspond directly to migrants' actual opportunities or outcomes. Examining how migrant characteristics shape more 
favorable or unfavorable beliefs, and how these beliefs relate to respondents' personal experiences and exposure to 
information, can provide evidence on how perceptions of different migrant groups are formed. The proposed design 
therefore provides a framework for understanding differences in beliefs about the integration of Ukrainian and Syrian 
migrants in Germany and, more broadly, the formation of beliefs about migrant integration.